\documentclass{article} % This command is used to set the type of
% document you are working on such as an article, book, or presenation
\usepackage{amssymb}
\usepackage{geometry} % This package allows the editing of the page layout
\usepackage{amsmath}  % This package allows the use of a large range
% of mathematical formula, commands, and symbols
\usepackage{graphicx}  % This package allows the importing of images

\newcommand{\question}[2][]{
  \begin{flushleft}
    \textbf{Question #1}: \textit{#2}

\end{flushleft}}
\newcommand{\sol}{\textbf{Solution}:} %Use if you want a boldface solution line
\newcommand{\maketitletwo}[2][]{
  \begin{center}
    \Large{\textbf{Assignment #1}

    Topology} % Name of course here
    \vspace{5pt}

    \normalsize{N. Ohayon Rozanes  % Your name here

    \today}        % Change to due date if preferred
    \vspace{15pt}

\end{center}}
\begin{document}
\maketitletwo[6]  % Optional argument is assignment number
%Keep a blank space between maketitletwo and \question[1]

\question[1]{}
\begin{enumerate}
  \item A space $X$ is Hausdorff if $\forall x, y  \in X, \exists$
    open neighborhoods $U_x, U_y$ such that $U_x \cap U_y =
    \emptyset$. Since open balls form a basis in a metric space, it
    suffices to show that $\exists r_x, r_y > 0$ such that
    $$
    B_{r_x} \cap B_{r_y} = \emptyset
    $$
    This can be achieved by first setting $r_x$ to be any positive
    real number less than $d(x,y)$, and setting $r_y = d(x,y) - r_x$.
    We confirm that the open balls centered at $x$ and $y$ with radii
    $r_x$ and $r_y$ are disjoint:
    \begin{align*}
      B_{r_x}(x) \cap B_{r_y}(y) &= \left\{z \in X | d(z, x) < r_x,
      d(z,y) < r_y\right\}\\
      &= \left\{z \in X | d(z, x) < r_x,
      d(z,y) < d(x, y) - r_x\right\}\\
      &= \left\{z \in X | d(z, x) < r_x,
      d(z,y) + r_x < d(x, y) \right\}\\
    \end{align*}
    Assume for the sake of contradiction that the above set is not
    empty, that is that there exists $z \in X$ such that
    \[
      d(z,y) + r_x < d(x,y) \quad d(z,x) < r_x
    \]
    Of course $d(z,x) < r_x \implies d(z,x) - r_x < 0$ then
    \[
      d(z,y) + r_x + d(z,x) - r_x < d(x,y) \implies d(z,y) + d(z,x) < d(x,y)
    \]
    However, by the triangle inequality
    \[
      d(x,y) \leq d(z,y) + d(z,x)
    \]
    Therefore, no such $z$ can exist, and the set
    \[
      B_{r_x}(x) \cap B_{r_y}(y) = \emptyset
    \]
    so any metric space $X$ is Hausdorff.
  \item A space is first countable if $\forall x \in X, x$ has a
    countable basis, that is, a countable collection of neighborhoods
    $\{ U_n \}$ such that if $U$ is a neighborhood of $x$, $\exists n$
    such that $U_n \subset U$.  Take the countable set of
    neighborhoods of a point $x\in X$
    \[
      \{U_n\} = \{ B_{\frac{1}{n}}(x)\}
    \]
    Since open balls form a basis in metric space, we have that the
    neighborhood $U$ of $x$ can be expressed as union of open balls
    of some radius $r$ centered on every point in $U$
    \[
      U = \bigcup_{y \in U} B_r(y)
    \]
    then there exists an open ball of some radius $r$ centered on $x$
    contained in $U$
    $$
    B_r(x) \subset U
    $$
    by the archimedian property of real numbers, $\exists n \in
    \mathbb{N}$ such that
    \[
      \frac{1}{n} < r \implies B_{\frac{1}{n}}(x) \subset B_r(x) \subset U
    \]
    Therefore, every point in a metrix space has a countable basis,
    and therefore every metric space is first countable.

\end{enumerate}
\question[2]{}
\begin{enumerate}
  \item
    \begin{enumerate}
      \item first we show that for every $x$, there exists $B \in
        \mathcal{B}$ such that $x \in B$. Take $\varepsilon > 0$ and $B =
        (x - \varepsilon, \infty)$ then $B \in \mathcal{B}$ and $x \in B$
      \item We wish to show that for any basis elements $B_a, B_b \in
        \mathcal{B}$ then $\exists B_c \in \mathcal{B}$ such that $B_c
        \subset B_a \cap B_b$. Let $a, b \in \mathbb{R}$ and $B_a = (a,
        \infty), B_b = (b, \infty)$, without loss of generality, assume
        that $a \leq b$. Then
        \[
          B_a \cap B_b = (b, \infty) = B_b
        \]
        Clearly then:
        \[
          B_b \subset B_a \cap B_b
        \]
        So take $B_c = B_b$. Therefore, $\mathcal{B}$ forms a basis as desired.
    \end{enumerate}
  \item Let $x, y \in X$, and $U_x, U_y$ be neighborhoods in the
    topology generated by $\mathcal{B}$. Then, since $U_x, U_y$ are
    open, there are open subset basis elements $(a, \infty) \subset
    U_x$ and $(b,\infty) \subset U_y$, then:
    \[
      (a, \infty) \cap (b,\infty) = (\min(a,b), \infty) \implies U_x
      \cap U_y \neq \emptyset
    \]
    Therefore, this space is not Haussdorff.
\end{enumerate}

\question[3]{}
\begin{enumerate}
  \item Let $U$ be an arbitrary open neighborhood of $0$, then
    $$
    U = \bigcup \prod U_i
    $$

    Where $U_i$ is either a basis element in $\mathbb{R}$ for a finite
    number of $i$ or a copy of $\mathbb{R}$.

    $$
    \prod U_i = (a_1, b_1) \times  (a_2, b_2) \times \hdots \times
    \mathbb{R} \times  \hdots
    $$

    Then for $0 \in U$, there must be at least one $0 \in \prod U_i
    \subset U$, that is:
    $$
    (0, 0, \hdots)\in(a_1, b_1) \times  (a_2, b_2) \times \hdots \times
    \mathbb{R} \times  \hdots
    $$

    As such every component of that product basis element must include
    $0$. $0 \in (a,b)\implies a < 0 < b$, as such $(a, b) \cap (0,
    \infty) = (0, b)$, therefore:
    \[
      \prod U_i \cap A = (0, b_1) \times (0, b_2) \times \hdots \times
      \mathbb{R}^+ \times \hdots \neq \emptyset
    \]

    Then, since $\prod U_i \subset U$:
    \[
      \prod U_i \cap A  \neq \emptyset \implies U \cap A \neq\emptyset
    \]

    As such, for any open neighborhood $U$ of $0$, the instersection of
    $U$ and $A$ is non-empty, and by definition $0$ is an
    accumulation point of $A$.
  \item A sequence in $A$ converging to $0$ in $\mathbb{R}^\omega$:
    $$
    a_n = \left(\frac{1}{n},\frac{1}{n}, \hdots\right)
    $$
    Clearly, $\forall n \in \mathbb{N}, \frac{1}{n} > 0 \implies a_n
    \in A$. Let $U$ be an arbitrary neighborhood of $0$, then, as in the last
    question, there exists a basis element subset $\prod U_i$ of $U$ such that
    \[
      \prod U_i = (a_1, b_1) \times (a_2, b_2) \times \hdots \times \mathbb{R}
      \times  \hdots
    \]
    Define $b$ as $\min(b_i)$, that is, the smallest right hand element of the
    finite number of finite intervals in the product basis. By the
    archimedian property of the real numbers, $\exists N \in
    \mathbb{N}$ such that $\frac{1}{N} < b$, then for any $n > N$:
    \[
      a_n \in  \prod U_i \subset U
    \]
    as such, $a_n \to 0$
\end{enumerate}
\end{document}
